\documentclass{article}
\usepackage{amsmath, amssymb, amsthm}
\usepackage{hyperref}
\usepackage{geometry}
\geometry{margin=1in}

\title{Erd\H{o}s Problem \#336}
\date{Last updated: 2025-08-31}
\author{Source: erdosproblems.com}

\begin{document}
\maketitle

\noindent\textbf{Status:} OPEN \quad \textbf{No prize}

\noindent\textbf{Tags:} number theory, additive basis

\medskip
\noindent\textbf{Problem Statement:}

For $r\geq 2$ let $h(r)$ be the maximal finite $k$ such that there exists a basis $A\subseteq \mathbb{N}$ of order $r$ (so every large integer is the sum of at most $r$ integers from $A$) and exact order $k$ (so every large integer is the sum of exactly $k$ integers from $A$). 

Find the value of\[\lim_r \frac{h(r)}{r^2}.\]

\bigskip
\noindent\textbf{Remarks:}

A simple example of the order of a basis differing from the exact order is given by $A=\cup_{k\geq 0}(2^{2k},2^{2k+1}]$, which has order $2$ but exact order $3$.

Erd\H{o}s and Graham \cite{ErGr80b} have shown that a basis $A$ has an exact order if and only if $a_2-a_1,a_3-a_2,a_4-a_3,\ldots$ are coprime. They also proved that\[\frac{1}{4}\leq \lim_r \frac{h(r)}{r^2}\leq \frac{5}{4}.\]The best bounds known for the limit are\[\frac{1}{3}\leq \lim_r \frac{h(r)}{r^2}\leq \frac{1}{2},\]the lower bound originally due to Grekos \cite{Gr88} and the upper bound to Nash \cite{Na93}. Improved bounds in the lower order terms were given by Plagne \cite{Pl04}.

Erd\H{o}s and Graham \cite{ErGr80b} showed $h(2)=4$. Nash \cite{Na93} showed $h(3)=7$. The value of $h(4)$ is unknown, but it is known \cite{Pl04} that $10\leq h(4)\leq 11$.

\bigskip
\noindent\textbf{References:}

[ErGr80b] Erd\H{o}s, P. and Graham, R. L., On bases with an exact order. Acta Arith. (1980), 201-207.
 
 [Gr88] Grekos, Georges, Sur l'ordre d'une base additive. ([1988?]), Exp. No. 31, 13.
 
 [Na93] Nash, John C. M., Some applications of a theorem of {M}. {K}neser. J. Number Theory (1993), 1--8.
 
 [Pl04] Plagne, Alain, \`A{} propos de la fonction {$X$} d'{E}rd\H{o}s et {G}raham. Ann. Inst. Fourier (Grenoble) (2004), 1717--1767.

\bigskip
\noindent\small{Source: \url{https://www.erdosproblems.com/336}}
\end{document}
